% Supplied Appendix E.4. Matrix target and allocation must be held consistent.
\subsection{Tail-Transverse Influence}
\label{app:tail-transverse-influence}
\label{sec:targeted-interventions}

An acquisition score evaluates how a candidate changes the current information audit. Fix the predictor, coordinates, calibration, slice rules, and finite audit weights $\pi_g\ge0$ summing to one. For action $a$, specify an admissible positive-semidefinite matrix path
\begin{equation}
J_g(t,a)=J_g+tH_{g,a}+o(t),\qquad t\downarrow0,
\label{eq:influence-matrix-path}
\end{equation}
where the remainder is in operator norm and $H_{g,a}$ is symmetric. The path incorporates allocation, normalization, and any retained pooling in the declared matrix target. Audit weights remain fixed along this path.

\paragraph{Schur differential and the lifted direction.}
Write
\begin{equation}
J_g=\begin{pmatrix}A_g&B_g\\B_g^\top&C_g\end{pmatrix},
\quad C_g\succ0,
\qquad
H=\begin{pmatrix}H_{TT}&H_{TN}\\H_{TN}^\top&H_{NN}\end{pmatrix}.
\label{eq:influence-block-partition}
\end{equation}
Differentiating the inverse in the Schur operator of Appendix~\ref{app:ridge-unregularized} gives
\begin{equation}
\begin{aligned}
DS_0(J_g)[H]
={}&H_{TT}-H_{TN}C_g^{-1}B_g^\top
-B_gC_g^{-1}H_{TN}^\top\\
&+B_gC_g^{-1}H_{NN}C_g^{-1}B_g^\top.
\end{aligned}
\label{eq:influence-schur-differential}
\end{equation}
Suppose $\lambda_g=\lambda_{\min}(S_0(J_g))$ is simple, with unit eigenvector $u_g$. Define the lifted efficient direction
\begin{equation}
p_g=\begin{pmatrix}u_g\\-C_g^{-1}B_g^\top u_g\end{pmatrix}.
\label{eq:influence-efficient-lift}
\end{equation}
Differentiating the eigenvalue equation and using $u_g^\top u_g=1$ gives $D\lambda_g[H]=u_g^\top DS_0(J_g)[H]u_g$. Expanding $p_g^\top Hp_g$ yields the four terms in this quadratic form. Therefore,
\begin{equation}
\boxed{
h_{g,a}:=
\left.\frac{d}{dt}\lambda_{\min}(S_0(J_g(t,a)))\right|_{t=0+}
=p_g^\top H_{g,a}p_g.
}
\label{eq:directional-schur-derivative}
\end{equation}
For an additive rank-one contribution $H_{g,a}=s_as_a^\top$,
\begin{equation}
\begin{aligned}
h_{g,a}&=(p_g^\top s_a)^2,\\
p_g^\top s_a
&=u_g^\top(s_{T,a}-B_gC_g^{-1}s_{N,a}).
\end{aligned}
\label{eq:influence-rank-one-gain}
\end{equation}
This measures the candidate score remaining along the weak task direction after nuisance residualization.

\paragraph{Allocation and normalization.}
For normalized augmentation by candidate matrix $M_{g,a}$,
\begin{equation}
J_g(t,a)=\frac{J_g+tM_{g,a}}{1+t},
\qquad H_{g,a}=M_{g,a}-J_g.
\label{eq:influence-normalized-path}
\end{equation}
The variational identity~\eqref{eq:reserve-schur-variational} gives $p_g^\top J_gp_g=\lambda_g$, so
\begin{equation}
h_{g,a}=p_g^\top M_{g,a}p_g-\lambda_g.
\label{eq:influence-normalized-gain}
\end{equation}
A replacement has perturbation equal to its new contribution minus its removed contribution, with the declared allocation rate included. These expressions place candidate gains in the same information units as the audited matrices.

\paragraph{Fixed-ridge derivative.}
For $S_{\rid}$ from Appendix~\ref{app:ridge-unregularized}, hold $\lambda_{\rid}$ fixed and set
\begin{equation}
Q_g=(C_g+\lambda_{\rid}I)^{-1},
\qquad
p_{g,\rid}=\begin{pmatrix}
u_{g,\rid}\\-Q_gB_g^\top u_{g,\rid}
\end{pmatrix},
\label{eq:influence-ridge-lift}
\end{equation}
where $u_{g,\rid}$ is a unit eigenvector associated with a simple smallest eigenvalue $\lambda_{g,\rid}$ of $S_{\rid}(J_g)$. Since $DQ_g[H_{NN}]=-Q_gH_{NN}Q_g$, the Schur differential takes the same form with $C_g^{-1}$ replaced by $Q_g$. Thus
\begin{equation}
D\lambda_{g,\rid}[H]=p_{g,\rid}^\top Hp_{g,\rid}.
\label{eq:influence-ridge-derivative}
\end{equation}
For normalized augmentation, the identity $Q_gC_gQ_g=Q_g-\lambda_{\rid}Q_g^2$ yields
\begin{equation}
p_{g,\rid}^\top J_gp_{g,\rid}
=\lambda_{g,\rid}
-\lambda_{\rid}\lVert Q_gB_g^\top u_{g,\rid}\rVert_2^2.
\label{eq:influence-ridge-baseline}
\end{equation}
Consequently,
\begin{equation}
\begin{aligned}
D\lambda_{g,\rid}[M_{g,a}-J_g]
={}&p_{g,\rid}^\top M_{g,a}p_{g,\rid}
-\lambda_{g,\rid}\\
&+\lambda_{\rid}\lVert Q_gB_g^\top u_{g,\rid}\rVert_2^2.
\end{aligned}
\label{eq:influence-ridge-normalized-gain}
\end{equation}
The ridge gain therefore uses its own eigenvalue, eigenvector, lifted direction, and allocation derivative consistently.

\paragraph{CVaR sensitivity.}
Choose either the exact or fixed-ridge matrix target consistently, and write its eigenvalues and derivatives as $\lambda_g$ and $h_{g,a}$. Fix $\varepsilon>0$ and $0<\alpha<1$. With $L_g=(\lambda_g+\varepsilon)^{-1}$, define the feasible tail weights
\begin{equation}
\mathcal W_{\alpha,\pi}
=\left\{\omega\ge0:\ \sum_g\omega_g=1,
\quad\omega_g\le\frac{\pi_g}{1-\alpha}\right\}.
\label{eq:influence-tail-weight-set}
\end{equation}
The finite-audit representation~\eqref{eq:app-finite-cvar} gives
\begin{equation}
C_\alpha(L):=\CVaR_{\alpha,\pi}(L_g)
=\max_{\omega\in\mathcal W_{\alpha,\pi}}\sum_g\omega_gL_g.
\label{eq:influence-cvar-dual}
\end{equation}
Let $\Omega_\alpha^\star(L)$ denote its maximizing weight vectors. If the maximizing vector $\omega_\alpha$ is unique and the smallest eigenvalues are simple, differentiation yields
\begin{equation}
\boxed{
G_\alpha(a):=-D_aC_\alpha
=\sum_g\omega_{\alpha,g}
\frac{h_{g,a}}{(\lambda_g+\varepsilon)^2}.
}
\label{eq:tail-information-gain}
\end{equation}
Indeed, $D_aL_g=-h_{g,a}/(\lambda_g+\varepsilon)^2$, and the unique active linear functional supplies the derivative of the CVaR maximization. The normalized weights include the upper-tail mass factor. A positive gain decreases inverse-information CVaR and increases its reciprocal:
\begin{equation}
D_aI_{\alpha,\pi}=I_{\alpha,\pi}^{\,2}G_\alpha(a),
\qquad I_{\alpha,\pi}=C_\alpha(L)^{-1}.
\label{eq:influence-certificate-derivative}
\end{equation}

\paragraph{Ties and repeated eigenvalues.}
At a repeated minimum eigenvalue, the exact right derivative uses its eigenspace $\mathcal E_g$:
\begin{equation}
h_{g,a}
=\min_{\substack{u\in\mathcal E_g\\\lVert u\rVert_2=1}}
p_g(u)^\top H_{g,a}p_g(u),
\label{eq:influence-repeated-eigenvalue}
\end{equation}
where the lift corresponds to the chosen exact or ridge target. This is the smallest eigenvalue of the Schur differential restricted to $\mathcal E_g$.

At a CVaR tie, the right derivative is obtained from the currently active linear functionals:
\begin{equation}
-D_a^+C_\alpha
=\min_{\omega\in\Omega_\alpha^\star(L)}
\sum_g\omega_g\frac{h_{g,a}}{(\lambda_g+\varepsilon)^2}.
\label{eq:influence-tied-cvar-gain}
\end{equation}
This follows because the right derivative of the maximum in~\eqref{eq:influence-cvar-dual} is the maximum of its active directional loss derivatives; negation turns that maximum into a minimum. Equations~\eqref{eq:influence-repeated-eigenvalue} and~\eqref{eq:influence-tied-cvar-gain} specify the directional gain at either source of nondifferentiability, including when both occur. A fixed maximizing tail-weight choice also defines an acquisition score; equality to the directional decrease follows under the unique-weight conditions above.

\paragraph{Tail-transverse influence.}
For a candidate intervention with assigned task label $y_a$, define
\begin{equation}
\boxed{
\TTI(a)=G_\alpha(a)\,\loss_w(x_a^{\mathrm{cf}},y_a),
}
\qquad\loss_w(x,y)=-\log p_w(y\mid x).
\label{eq:tail-transverse-influence}
\end{equation}
The information factor emphasizes weak task directions in the current audit. The loss factor emphasizes candidates that the current predictor handles poorly. Their product ranks candidates under the declared acquisition budget. At nonsmooth points, the directional gain in~\eqref{eq:influence-tied-cvar-gain} supplies an exact right-sensitivity convention; an implementation using a fixed tail choice records that choice.

Scores are evaluated at the fixed predictor, and selected candidates then enter training. Held-out evaluation measures the realized information and predictive-performance changes after finite augmentation and retraining. The reported feature-space experiment uses a normalized residual-square surrogate $\widetilde G_\alpha(a)$ in its product score. Appendix~\ref{app:feature-acquisition} specifies that implemented score and its distinction from the directional CVaR decrease derived here.
